\section{Isospectrality without bounded similarity}
\label{sec:similarity}

The spectral idempotent for the nonzero eigenvalue in block $m$ is
\begin{equation}
  \Pi_{T,m}=\lambda_m^{-1}T_m.
  \label{eq:riesz-block}
\end{equation}
It is an oblique projection onto $\C e_m$, and therefore
\begin{equation}
  \|\Pi_{T,m}\|=\frac{\rho_T(m)}{|\lambda_m|}.
  \label{eq:riesz-ratio}
\end{equation}
Combining this ratio with the two block masses yields
\begin{align}
  \|\Pi_{S,m}\|
    &=\prod_{p\in J_h(m)}(1-p^{-\sigma})^{-1/2},
      \label{eq:riesz-S}\\
  \|\Pi_{M,m}\|&=\sqrt{\zeta(h\sigma)}.
      \label{eq:riesz-M}
\end{align}

The following block lemma supplies the converse that a mere appeal to
necessary projection bounds would miss.

\begin{lemma}[Uniform block diagonalization]
\label{lem:uniform-diagonalization}
Let $T=\bigoplus_mT_m$ be a compact orthogonal direct sum of rank-one
blocks.  Suppose each block has nonzero eigenvalue $\lambda_m$ with range
$\C e_m$, and suppose that the $\lambda_m$ are pairwise distinct.  Then
$T$ is boundedly similar to the normal diagonal operator
with entries $\lambda_m$ and zeros on the block kernels if and only if
$\sup_m\|\lambda_m^{-1}T_m\|<\infty$.
\end{lemma}

\begin{proof}
If $T=XNX^{-1}$ and $N$ is normal, the nonzero spectral projections of
$N$ are orthogonal.  Hence those of $T$ have norms at most
$\|X\|\|X^{-1}\|$.

Conversely, decompose a block as
$\Hs_m=\C e_m\oplus K_m$.  It has the matrix form
\[
  T_m=\begin{pmatrix}\lambda_m&\varphi_m\\0&0\end{pmatrix}.
\]
Set
\[
  Y_m=\begin{pmatrix}1&-\lambda_m^{-1}\varphi_m\\0&I_{K_m}\end{pmatrix}.
\]
A direct multiplication gives
$Y_m^{-1}T_mY_m=\diag(\lambda_m,0)$.  The norm of
$\lambda_m^{-1}\varphi_m$ is controlled by the norm of the idempotent
$\lambda_m^{-1}T_m$; therefore the hypothesis uniformly bounds both
$Y_m$ and $Y_m^{-1}$.  The direct sum $Y=\bigoplus_mY_m$ is bounded and
invertible and conjugates $T$ to the stated diagonal.  Since
$|\lambda_m|\to0$, that diagonal is compact and normal.  A norm identity
and explicit bounds are given in \cref{app:graph-transform}.
\end{proof}

\begin{theorem}[Exact similarity gap]
\label{thm:similarity}
\begin{align*}
  S_{h,s}\sim_{\mathrm{bd}}\text{ compact normal}&\iff\sigma>1,\\
  M_{h,s}\sim_{\mathrm{bd}}\text{ compact normal}&\iff\sigma>1/h.
\end{align*}
\end{theorem}

\begin{proof}
The modulo norm in \cref{eq:riesz-M} is a finite constant, independent of
$m$, throughout the bounded domain.  The lemma applies.

Every finite set of primes can occur as $J_h(m)$, by taking
$m=\prod_{p\in E}p^{h-1}$.  Consequently
\[
  \sup_{m\in\Fh}\|\Pi_{S,m}\|
  =\left(\prod_p(1-p^{-\sigma})^{-1}\right)^{1/2}
\]
as a finite value or $+\infty$.  This is finite exactly for $\sigma>1$.
The lemma again gives both directions.
\end{proof}

\begin{figure}[t]
  \centering
  \resizebox{0.98\linewidth}{!}{\begin{tikzpicture}[x=2.15cm,y=0.82cm,>=Stealth]
  \draw[->] (-0.12,0) -- (3.25,0) node[right] {$\sigma$};
  \foreach \x/\lab in {0/0,1/{1/h},2/1}
    \draw (\x,0.10)--(\x,-0.10) node[below=2pt] {$\lab$};
  \draw[densely dotted] (1,0)--(1,4.55);
  \draw[densely dotted] (2,0)--(2,4.55);
  \node[anchor=east] at (-0.12,1) {$S$ bounded};
  \draw[very thick,blue!65!black] (0.02,1)--(3.05,1);
  \node[anchor=east] at (-0.12,2) {$M$ bounded};
  \draw[very thick,orange!80!black] (1.02,2)--(3.05,2);
  \node[anchor=east] at (-0.12,3) {$S\sim N$};
  \draw[very thick,blue!65!black] (2.02,3)--(3.05,3);
  \node[anchor=east] at (-0.12,4) {$M\sim N$};
  \draw[very thick,orange!80!black] (1.02,4)--(3.05,4);
  \node[align=center,fill=red!6,draw=red!55!black,rounded corners]
    at (1.5,5.2) {$1/h<\sigma\le1$: same nonzero spectrum,\\$M$ similar to normal, $S$ not similar};

  \node[draw=black!55,rounded corners,align=left,anchor=north west,
        font=\small,text width=80mm,minimum height=11mm] at (-0.1,-1.00)
    {Power ideals:\\$S^k\in\mathcal S_q\iff k\sigma q>2$\\
     $M^k\in\mathcal S_q\iff \sigma>1/h$ and $k\sigma q>2$};
  \node[draw=black!55,rounded corners,align=left,anchor=north west,
        font=\small,text width=80mm,minimum height=11mm] at (-0.1,-2.75)
    {Commutator ideals:\\$[S^*,S]\in\mathcal S_q\iff\sigma q>1$\\
     $[M^*,M]\in\mathcal S_q\iff \sigma>1/h$ and $\sigma q>1$};
\end{tikzpicture}
}
  \caption{Exact existence, similarity, power-ideal, and commutator-ideal
  conditions.  The horizontal bands encode only the fixed ordering
  $0<1/h\le1$; the ideal thresholds are stated algebraically because their
  relative positions depend on $h,k,q$.  Every displayed wall is strict.}
  \label{fig:phase}
\end{figure}

In the full interval $1/h<\sigma\le1$, both operators exist and have the
same simple nonzero spectrum.  Their common power traces and determinants
also agree whenever the corresponding ideal condition is met; here
``determinants'' means order-$r$ regularized determinants for integers
$r$ satisfying $r\sigma>2$, never the order-one Fredholm determinant.  Yet one
operator has uniformly controlled eigenprojections and the other does not.
This is an obstruction to bounded similarity, not merely to unitary
equivalence.
